\documentclass[12pt]{article}
\usepackage[utf8]{inputenc}
\usepackage[T1]{fontenc}
\usepackage{hyperref}
\title{Tanapag Language Resources for K-12 Curriculum Development}
\author{Compiled by Claude Sonnet 5 (claude-sonnet-5) for Andr\'as Kornai}
\date{}
\begin{document}
\maketitle

\section{Introduction}
Tanapag (ISO 639-3 tpv; Glottocode tana1281) is an Austronesian language of Saipan, Northern Mariana Islands. It is also called Northern Carolinian or Talaabwogh. Sources disagree on whether it is a separate language or a variety of Carolinian, and on how many speakers it has. Omniglot reports about 10 speakers, while a Rosetta Project page gives about 4,400, which probably counts Carolinian speakers more broadly. Younger people tend to use Chamorro. Very little material exists specifically for Tanapag.

\section{Primary Educational Resources}
I found no Tanapag textbooks, primers or school programmes.

\section{Linguistic and Reference Materials}
\begin{itemize}
\item Jackson and Marck, \emph{Carolinian-English Dictionary} (PALI Language Texts, University of Hawaii Press). Wikipedia says it includes some Tanapag forms, but I did not check whether it is in print. It is cited on the Wikipedia page: \url{https://en.wikipedia.org/wiki/Tanapag_language}
\item Saipan Carolinian (Tanapag) wordlist in the Austronesian Basic Vocabulary Database: \url{https://abvd.eva.mpg.de/austronesian/language.php?id=339}
\item Rosetta Project page for tpv: \url{https://rosettapanglossia.longnow.org/index.php?title=Tpv}
\item Carolinian numbers, some in Tanapag (Omniglot): \url{https://www.omniglot.com/language/numbers/carolinian.htm}
\end{itemize}

\section{Religious and Community Resources}
The search found none.

\section{Digital and Audio Resources}
The Pacific Worlds CNMI site has Carolinian vocabulary tables, which are not labelled as Tanapag. I found no audio resources.
\begin{itemize}
\item \url{https://www.pacificworlds.com/cnmi/native/lang2.cfm}
\item \url{https://www.pacificworlds.com/cnmi/arrival/lang1.cfm}
\end{itemize}

\section{Curriculum Development Context}
Tanapag is nearly extinct, and I found no teaching materials for it. Any curriculum would need to start with documentation work with the remaining speakers. It could draw on the broader Carolinian materials and the Carolinian-English Dictionary. Whether Tanapag is treated as separate from Carolinian should be settled with the community.
\end{document}